\subsection{Artifacts a model actually wrote}
\label{sec:llm}

Both experiments above generate artifacts from a distribution we chose. This one
does not. Three agents review a real code change --- a connection-pool size
reduction shipped alongside an unrelated fix --- each taking a distinct aspect.
With $S_i = 1$ the agent is given the diff and must cite identifiers or line
numbers; with $S_i = 0$ it is told to produce confident, professional review
guidance \emph{without} the code and not to admit it. That is shirking in the
sense this paper means: the artifact still looks like work. A coordinator reads
the three reviews and ships or holds. The records are derived from the episode
exactly as in Section~\ref{sec:episodes}, a judge model predicts which agents
actually read the diff, and the paper's instruments are applied to the induced
channel. Model \texttt{codestral-2508} at temperature $0$; $192$ episodes, $24$
per effort profile; all artifacts, records and verdicts are in the artifact.

\begin{table}[t]
\centering
\caption{Attribution error on model-written reviews, $192$ episodes. The first
and third rows contain \emph{the same text}, differing only by ``Reviewer $k$:''
labels.}
\label{tab:llm}
\small
\begin{tabular}{@{}lrr@{}}
\toprule
record & bytes & attribution error (of 3 bits) \\
\midrule
full trace              & $2612$ & $0.391$ \\
per-agent summary       & $661$  & $0.260$ \\
artifact multiset       & $2582$ & $1.333$ \\
global summary          & $609$  & $1.000$ \\
final output only       & $507$  & $1.193$ \\
\bottomrule
\end{tabular}
\end{table}

\paragraph{Thirty bytes of authorship.}
The full trace and the artifact multiset hold the same five reviews. They differ
by about $30$ bytes of attribution labels, a $1.2\%$ difference in size, and by
$0.943$ bits of attribution error --- $95\%$ bootstrap interval
$[0.802, 1.089]$ over $20{,}000$ resamples, excluding zero. Roughly a third of
the recoverable per-agent signal lives in the labels rather than in the content,
on text a language model wrote. This is the identity-versus-volume finding of
Section~\ref{sec:episodes} reproduced outside the generative model that produced
it, and it is the claim this experiment is offered to support.

\paragraph{What we do not claim from it, and why.}
The measured margins do \emph{not} order as the synthetic experiments do: the
full trace comes second from bottom. Decomposing $\gamma$ by agent shows why, and
the cause is our task rather than the records. The worst single-cell total
variation is $0.958$ for the first agent and $0.792$ for the second, but $0.125$
for the third, whose sub-question asks what is now untested --- a question whose
content-free answer is nearly indistinguishable from a real one. Because
$\gamma$ is a \emph{minimum} over agents, that one sub-question sets the margin
for every record, and the ordering across records is then a statement about our
question set, not about compression.

We record this as a methodological constraint on measuring $\gamma$ empirically
at all: an all-teammate-profile margin is only as informative as the
\emph{least auditable} sub-task in the episode, so sub-tasks must be matched for
auditability before any measured $\gamma$ can be attributed to the record.
Attribution error, which averages over agents rather than minimising, does not
have this defect, which is why Table~\ref{tab:llm} reports it.

\paragraph{One honest artefact of sample size.}
At four episodes per profile the same estimator returned
$\gamma(\text{full trace}) = 0$, triggering the failure condition registered
before the run. The cause was seven cells in which the judge happened to give
identical verdicts with and without some agent's effort, forcing an exact zero;
such ties can only bias $\gamma$ downward. They fell to one cell at ten episodes
per profile and to none at twenty-four. The attribution gap, by contrast, moved
from $0.969$ to $0.943$ across that sixfold increase in data.
